\subsection{A bounded exponent on the actual primitive factor}
\label{sec:bounded-primitive-exponent}

\paragraph{Finding A47 (bounded primitive torsion and its exact obstruction).}
Let $C/k$ be a smooth projective geometrically connected curve
of genus $g\ge2$ over a finite field of characteristic $p$.
Let $P\in C(k)$, and let $D\to C$ be a geometrically connected
finite etale cover of odd degree $n>1$ prime to $p$, geometrically
Galois with cyclic group. Assume that $P'\in D(k)$ lies above $P$.
The geometric deck group need not be constant over $k$.
Let $P_n\subset J_D$ be the actual primitive conductor-$n$
subvariety defined by the integral projector (IS.2)--(IS.3), and put
\[
 s=\omega(n),\qquad c=2^s,\qquad
 M_n=\exp\bigl(P_n(k)[p^\infty]\bigr),
 \tag{BP.1}
\]
with exponent one for the zero group. Suppose that
\[
                         n>(cM_n-1)^2.
 \tag{BP.2}
\]
All section statements below concern the geometrically
pro-prime-to-$p$ arithmetic fundamental group.
Then a collision between the point sections of $P$ and a distinct
$Q\in C(k)$ forces a separable map to an elliptic curve
\[
 C_{\bar k}\longrightarrow E,
 \qquad \deg(C_{\bar k}/E)=p^a=e>1,
 \qquad e\mid M_n,\quad e\le g-1,
 \tag{BP.3}
\]
totally ramified at both $P$ and $Q$. Moreover the given geometric
cover $D_{\bar k}\to C_{\bar k}$ is the pullback of an etale
cyclic cover $Z\to E$ of degree $n$ between genus-one curves.
Consequently (BP.2) separates the section of $P$ from all other
rational point sections whenever $2\le g\le p$, or whenever
there is no map (BP.3) totally ramified at $P$.

If $M_n=1$, the conclusion is separation in every genus
$g\ge2$: (BP.3) is then impossible, and (BP.2) is automatic
by (IS.11). Thus the rational primitive criterion A46 is
recovered, including its nonconstant-deck scope. No existence
of a cover satisfying (BP.2) on every prescribed curve is asserted.

\begin{proof}
\emph{1. The principal projector divisor and its exact order.}
Suppose that the sections of $P$ and $Q\ne P$ collide.
The same pointed-cover and finite-field Kummer implication as
in A46 gives $Q'\in D(k)$ with
\[
                 \delta=[Q'-P']\in J_D(k)[p^\infty].
 \tag{BP.4}
\]
Choose a geometric generator $\sigma$ and put
\[
 A=\prod_{\ell\mid n}(1-\sigma^{n/\ell}),\qquad
 p_n=ne_n\in\mathbf Z[\langle\sigma\rangle],
 \tag{BP.5}
\]
where $e_n$ is the faithful-character projector. The actual
endomorphism $p_n$ is defined over $k$, and $Ap_n=nA$ by
(IS.2)--(IS.5). Hence
\[
 M_np_n\delta=0,\qquad nM_nA\delta=0,
                         \qquad M_nA\delta=0.
 \tag{BP.6}
\]
The last implication uses that $A\delta$ is $p$-primary and
$p\nmid n$. In particular the geometric operator $A$ is used
only after a rational bound has been obtained on the actual
subvariety over $k$; the field has not been enlarged.

Let $m$ be the exact order of $A\delta$, with $m=1$ for zero.
Then $m$ is a power of $p$ and $m\mid M_n$. The expansion of
$A$ has $c=2^s$ distinct coefficients, each $1$ or $-1$, half
of each sign, as in A46. The two deck orbits of $P'$ and $Q'$
are disjoint and free. There is therefore a function $f$ on
$D_{\bar k}$ with
\[
 \operatorname{div}(f)=mA(Q'-P'),\qquad \deg f=d=cm.
 \tag{BP.7}
\]
Every nonzero divisor coefficient is $m$ or $-m$.
The function is separable. For $m>1$, a $p$th power would
make $(m/p)A(Q'-P')$ principal, contradicting exact order;
for $m=1$, the simple poles exclude a $p$th power directly.

\emph{2. The controlled field and free cyclic action.}
Set
\[
 K=\bar k(D),\qquad
 F=\bar k(\gamma f:\gamma\in\langle\sigma\rangle),
                  \qquad e=[K:F].
\]
Let $Z$ be the smooth projective curve with field $F$ and
$v:D_{\bar k}\to Z$ the associated map. The field is generated
by finitely many degree-$d$ elements over the perfect field
$\bar k$. Poonen's Corollary~2.3 and its definition of
$d$-controlled therefore give exactly
\[
                 e\mid d,\qquad g(Z)\le(d/e-1)^2
 \tag{BP.8}
\]
\cite[Proposition~2.1, Lemma~2.2 and Corollary~2.3]{Poonen2007Gonality}.
Since $g(D)=n(g-1)+1$ and (BP.2) implies $n>(d-1)^2$,
one has $e>1$. Moreover $e\mid2^sp^b$ for some $b$, so
$\gcd(e,n)=1$. The extension $K/F$ is separable, being
intermediate in $K/\bar k(f)$.

The cyclic group preserves $F$ and acts faithfully on $Z$:
the order of a kernel divides both $e$ and $n$.
Its action is free. If a point stabilizer contained an element
of prime order $\ell\mid n$, its free action on $D$ would
partition the fibre divisor above that point into orbits of
length $\ell$ with equal multiplicities on each orbit.
This would give $\ell\mid e$, a contradiction.
Consequently the quotient $E=Z/\langle\sigma\rangle$ is etale
of degree $n$. The bound $0\le g(Z)<n$ and
\[
                         g(Z)-1=n(g(E)-1)
\]
force $g(Z)=g(E)=1$. In particular
\[
                 \deg\operatorname{Diff}(v)=2n(g-1).
 \tag{BP.9}
\]

\emph{3. Integral signs exclude identical image orbits.}
The maps from each of the two original deck orbits into $Z$
are injective, because its cyclic action is free. Their images
are either identical or disjoint.
Suppose first that they are identical, and choose $j$ with
\[
                         v(P')=\sigma^jv(Q').
\]
Let $r_P,r_Q$ be the ramification indices of $v$ at $P',Q'$;
they are positive and constant along the respective deck orbits.
Write $f=h\circ v$. On the image orbit based at $v(Q')$, put
\[
 H=\sum_{i=0}^{n-1}
       \operatorname{ord}_{\sigma^iv(Q')}(h)\,\sigma^i
                         \in\mathbf Z[\langle\sigma\rangle].
\]
Comparing the literal integer coefficients of (BP.7) on the
two free orbits gives
\[
           mA=r_QH,\qquad -mA=r_P\sigma^{-j}H.
\]
Eliminating $H$ and cancelling the positive integer $m$ in
the torsion-free integral group ring yields
\[
                        r_QA=-r_PA\sigma^{-j}.
 \tag{BP.10}
\]
If $t$ is the order of $\sigma^j$, then $t$ is odd. Iteration
gives $r_Q^tA=-r_P^tA$, which is impossible: $r_Q^t+r_P^t$
is positive and $A\ne0$. This argument does not assume a
quadratic map, equal indices, or reduced fibres. Its signs
are integer divisor signs, also in characteristic two.

\emph{4. Disjoint image orbits force the wild elliptic obstruction.}
The images must therefore be disjoint. Take any point in the
support of (BP.7). A fibre above its image can contain only
points in that divisor's support, since it is a zero or pole
fibre of $h$. It contains at most one point from either original
orbit, by injectivity, and cannot contain one from the other
orbit, by disjointness. It is consequently a singleton fibre
of total ramification index $e$. Its divisor multiplicity
$m$ is a multiple of $e$, so
\[
                         e\mid m\mid M_n.
 \tag{BP.11}
\]
Since $e>1$, it is a nontrivial power of $p$.
The identity coefficient of $A$ is one; in particular $P'$
and $Q'$ themselves belong to the support. Equivariance makes
every point of their two disjoint orbits totally ramified
with index $e$.

At such a point, a local parameter expression
$u=a_et^e+\cdots$ has nonzero derivative by separability,
but its degree-$e$ leading derivative vanishes since $p\mid e$.
Thus the different exponent is at least $e$; this is the local
different formula in the Riemann--Hurwitz derivation
\cite[Lemma~53.12.4, Tag~\texttt{0C1F}]{StacksRH}.
The two orbits have $2n$ points in total, and (BP.9) gives
\[
                         2ne\le2n(g-1),
                    \qquad p\le e\le g-1.
 \tag{BP.12}
\]
This is impossible for $2\le g\le p$, and is also impossible
when $M_n=1$ by (BP.11).

\emph{5. The obstruction descends to the given base curve.}
Let $G=\langle\sigma\rangle$.
The field $K/K^G$ is Galois with group $G$. Its subgroup fixing
$FK^G$ is the kernel of the faithful action on $F$, so
\[
             K=FK^G,\qquad \bar k(C)=K^G,
                 \qquad \bar k(E)=F^G.
 \tag{BP.13}
\]
It follows that $[\bar k(C):\bar k(E)]=e$. This extension is
separable, since $K/F$ and $F/F^G$ are separable. Choosing an
origin on the genus-one curve $E$ makes it an elliptic curve.
The etale fibre product $C_{\bar k}\times_E Z$ is connected
and has function field $K$, hence is $D_{\bar k}$.
Ramification indices are unchanged by the etale base change
$Z\to E$. The total ramification at $P',Q'$ therefore gives
total ramification at $P,Q$ in the degree-$e$ map $C_{\bar k}\to E$.
This proves every assertion in (BP.3) and the theorem.
\end{proof}

\paragraph{Why the remaining configuration is not a sign contradiction.}
On an elliptic curve a fixed-point-free finite cyclic action
consists of translations. Indeed an automorphism with a nontrivial
origin-preserving part has a fixed point, since one minus that
part is a surjective nonzero endomorphism. Translations act
trivially on $\operatorname{Pic}^0$. Thus, in the configuration
of (BP.3), $A(y-x)$ is principal for any two points $x,y$ of
$Z$. If $x=v(P')$ and $y=v(Q')$, total ramification gives
\[
                   v^*A(y-x)=eA(Q'-P').
 \tag{BP.14}
\]
The remaining positive-power principal divisor is therefore
compatible with this elliptic pullback geometry. This explains
why the proof retains (BP.3) when $M_n>1$; it does not assert
that such a divisor arises from a full point-section collision,
or that an arbitrary elliptic pullback satisfies (BP.2).

\paragraph{The precise remaining witness requirement.}
For the characteristic-two genus-two part of U01 and
\texttt{Q-HOOLEY}, including the ordinary real-multiplication
targets of A33, A46's sufficient input $M_n=1$ can now be
replaced by a pointed cover satisfying
\[
                 M_n<\frac{\sqrt n+1}{2^{\omega(n)}}.
 \tag{BP.15}
\]
The primitive factor may have positive geometric two-rank and
nonzero rational two-primary torsion. This changes the sufficient
input to the point-section route; it does not solve the literal
primitive-rank-zero existence problem or the rational-unit
witness questions in the recovered A1 rows B.24, B.5 and B.32.

Compared with A45, the whole-Prym exponent $M$ has been replaced
by $M_n$, at the cost of replacing $2$ by $2^{\omega(n)}$ in
the bound. The primitive test has a smaller threshold when
$M>2^{\omega(n)-1}M_n$; in other cases the two numerical tests
need not imply one another. For prime-power degree the coefficient
is still two, and only the top primitive conductor is measured.
No general estimate supplying (BP.15) for every prescribed curve
or point follows from this proof or from Hooley's prime-density
theorem. In particular, arbitrary constant extension is not a
way to improve a rational torsion exponent.
